\odiseachapter{Decimocuarta parte}{Penélope}

El penúltimo canto de la \emph{Odisea} arranca con Euriclea dándole a Penélope la noticia de que su marido está de vuelta:

\begin{verse}
«Despierta, Penelopea, mi niña, para que veas\\
tú con tus ojos aquello que todos los días anhelas.\\
Vino Odiseo, ha llegado a su casa, aunque tarde volviera,\\
y a los altivos galanos mató, los que casa y hacienda\\
roían y devoraban, vejando a tu hijo sin tregua».
\end{verse}

La reacción de Penélope es comprensible. Después de veinte años de espera, ya no se hace ilusiones. Pero todavía más interesante es su opinión sobre los dioses, a los que no tiene en muy alta estima:

\begin{verse}
Y le responde en un verbo Penélope milpensamientos:\\
Nana querida, te habrán vuelto loca los dioses, que a ellos\\
nunca les falta el poder de volver insensato al más cuerdo\\
y hacer entrar en cordura al flojo de entendimiento;\\
daño te hicieron, que a ti antes el seso te andaba a derecho.
\end{verse}

Después de declarar su poca fe en los olímpicos, intenta despedir a la nodriza con una suave riña:

\begin{verse}
¿Por qué te burlas de mí, que tengo el alma en un duelo,\\
con ese hablar de través sacándome a mí de este sueño\\
dulce que me encadenaba cubriendo mis párpados buenos?\\
Que nunca había dormido así desde que Odiseo\\
a ver marchó Malatroya, que ni nombrar esa quiero.\\
Anda, ve abajo ahora mismo y torna a la sala de nuevo,\\
pues que de haber sido otra de las mujeres que tengo\\
la que con este recado me hubiera sacado del sueño,\\
a fe que a las duras la habría mandado volver en despecho\\
a sus faenas: en algo vejez te sirvió de provecho.
\end{verse}

Euriclea le explica a su ama que Odiseo no es otro que el mendigo que ella ha favorecido un rato antes. Penélope le pide detalles de la batalla y la dulce ancianita, que no tiene pelos en la lengua, se los ofrece:

\begin{verse}
Hallé a Odiseo en mitad de los cuerpos que él mismo abatiera,\\
erguido, y a su redor, abrazando la muy dura tierra,\\
unos sobre otros, los muertos ---¡qué gozo esa vista te diera!---\\
y embadurnado él de sangre y matanza tal si un león fuera.\\
Ante las puertas del patio ahora todos se encuentran\\
amontonados; sahumó él con azufre la estancia rebella\\
luego de hacer un gran fuego, y aquí me mandó que viniera.
\end{verse}

Pero Penélope sigue sin creer que Odiseo haya vuelto y prefiere atribuir la matanza a algún dios airado, con tal de no ceder a la esperanza.

\begin{verse}
no creo que esto que cuentas pasara en verdad de esa guisa,\\
que fue sin duda algún dios quien mató a la galana cuadrilla\\
en ira por su maldad y su cordipunzante osadía.\\
Respeto a hombre ninguno de los de la tierra tenían\\
por vil o noble que fuera el que a su mesa venía;\\
por eso ellos penaron en mal, por su altanería;\\
y allende Acaya Odiseo el regreso perdió, y aun la vida.
\end{verse}

En la \emph{Odisea} original no hay colgante que permita que los esposos se reconozcan, pero sí aparece la célebre cicatriz dejada por el jabalí. Euriclea le dice a Penélope:

\begin{verse}
Mas, ¡ea!, voy a decirte una seña que es bien manifiesta:\\
la cicatriz que del blanco colmillo de un puerco le queda,\\
que la noté yo al lavarle sus pies y a decírtelo fuera\\
cuando cerró él con su mano mi boca, que no consintieran\\
que te dijera yo nada sus mientes biensabideras.
\end{verse}

Ese detalle la ablanda, pero sigue sin dar su brazo a torcer. Acepta, sin embargo, bajar al escenario de la carnicería. El encuentro entre los esposos sorprende por la compostura, casi la timidez que ambos manifiestan:

\begin{verse}
Cuando llegó y los lumbrales de piedra pasó, ya en la sala,\\
sentose frente a Odiseo, al resplandor de las llamas,\\
pegada a la otra pared; contra un largo pilar él estaba\\
sentado, al suelo mirando, a la espera de que una palabra\\
su alta esposa dijera cuando a sus ojos se hallara.\\
Largo duró ella sin voz, la extrañeza roía su entraña;\\
y a veces era a Odiseo a quien veía en su cara\\
y a veces ya no era él por el argamandel que llevaba.
\end{verse}

El hijo pierde la paciencia con la madre y la increpa así:

\begin{verse}
«Madre mía, madrastra, que áspero espíritu tienes,\\
¿por qué tan lejos del padre te sientas en vez de ponerte\\
cerca de él a escuchar lo que quieres saber y requieres?\\
Ninguna otra mujer se tendría con ánimo inerte\\
sin acercarse al esposo que, luego de mil padeceres,\\
vuelve a su patria y su tierra tras veinte años ausente.\\
Más duro siempre que piedra el corazón que tú tienes».
\end{verse}

A lo que Penélope responde:

\begin{verse}
«Mi corazón, hijo mío, está como en pasmo en mi pecho,\\
y ni decirle palabra ni hacerle pregunta ahora puedo,\\
ni tan siquiera mirarlo a la cara de frente. Si es cierto\\
que es Odiseo, que ha vuelto a su casa, nosotros dos luego\\
reconocernos podremos entrambos mejor, pues tenemos\\
señas que son solo nuestras y ocultas están para el resto».
\end{verse}

Nada más lejos de la sacarina que le echa Nolan al amargo café de este reencuentro. Nada de saltar a los brazos del marido, nada de encendidos besos. Penélope no se atreve casi ni a mirarle a los ojos. ¿Es de extrañar? Lleva dos décadas sola, criando a su hijo como puede, sorteando con mil artes el acoso de los ahora difuntos pretendientes. Hace años que los últimos aqueos han regresado y es totalmente comprensible que no albergara ya esperanza alguna de reencontrarse con un hombre que lleva ausente de casa la mitad de su vida. Penélope quizás tendría catorce o quince años cuando se casó con Odiseo, puede que tuviera dieciséis cuando nació Telémaco y este era todavía un bebé cuando Ulises se marchó. ¿Cuánto tiempo vivió con él, dos años? ¿Cómo se comparan con las dos décadas de ausencia?

¿Cuán ferviente es el deseo de Penélope de recuperar a su esposo? Ciertamente no tenía intención de desposar a uno de los galanes, pero eso lo podemos comprender perfectamente, dada la mala catadura de casi todos ellos. E incluso si hubiera alguno pasable, ¿para qué? ¿Qué falta le hace otro hombre en su vida, cuando ya tiene a su hijo? Quizás Penélope, la que teje y desteje un sudario, lo que desea es seguir esperando, es decir, que la dejen en paz.

Pero no ha sido ese el deseo de los dioses, o eso parece. Penélope sigue en sus trece y dispone de recursos para averiguar, en privado, si es de verdad su marido el que ha vuelto, o se trata de un divinal impostor. Odiseo, que sabe de qué prueba se trata, calma al hijo:

\begin{verse}
«Deja que sea tu madre la que en el palacio, Telémaco,\\
me ponga a prueba; que ya vendrá ella a mejor pensamiento.\\
Ahora voy sucio y me cubro con míseros trapos el cuerpo,\\
y aún no cree que soy él, pues me tiene por hombre de menos.\\
Pensemos mientras tú y yo qué buen fin puede darse a lo nuestro.\\
Pues si hasta aquel que ha matado a un hombre no más en el pueblo\\
---a uno que a muchos no deja tras él que le venguen el duelo---\\
huye de tierra y de patria, y desampara a sus deudos,\\
¿qué de nosotros, que a flor y bastión de la villa hemos muerto\\
de Ítaca belmoceril? Yo te llamo a que pienses en esto».
\end{verse}

El metisagaz Ulises está, de hecho, más preocupado por cómo se va a tomar el pueblo la matanza de un centenar de príncipes ---aunque se da el caso de que muchos de ellos son vástagos de casas nobles de las islas vecinas--- que por quedarse a solas con su esposa. La salida del padre tiene la virtud de centrar a Telémaco, quien dice:

\begin{verse}
«Padre querido, eso míralo tú, pues que más astucioso\\
que tu saber entre hombres no hay, y mortal no habría otro\\
que aventurárase a entrar a porfía contigo. Nosotros\\
detrás iremos de firme, y en nada, yo creo, el arrojo\\
ahí nos ha de faltar, mientras haya de fuerzas un soplo».
\end{verse}

Como ya te imaginas, resabida lectora, vueltodetodo lector, milardido lo tiene todo previsto:

\begin{verse}
«Pues bien, yo voy a decir lo que a mí me parece más cuerdo.\\
Para empezar os bañáis y os aparejáis jubón nuevo,\\
luego mandáis a las criadas de casa a que escojan atuendo,\\
y que a su cítara clarisonera el aedo del cielo\\
nos guíe a todos en danza bien viva y amiga de juego:\\
que se figure una boda quien sea que esté fuera oyendo,\\
ya sea gente de paso o los que viven linderos.\\
Que no ande en bocas por la ciudad la matanza de estos,\\
los pretendientes, hasta que a nosotros nos dé abrigadero\\
nuestra masada milarbolada; y allí ya veremos\\
qué otra ventaja va a dar a la mano el Olímpico luego».
\end{verse}

Dicho y hecho, los hombres se lavan y se ponen ropa limpia, las sirvientas se acicalan ---quizás a más de una se le corriera el rímel por culpa de las lágrimas que se le escaparon al pasar junto al corral donde colgaban las doce infortunadas---, el aedo echa mano de la cítara y en menos de lo que se degüella a un inocente, han montado un estupendo jolgorio. Los que oyen la algarabía, desde fuera, creen saber lo que está pasando:

\begin{verse}
«A fe que alguno ha casado por fin con la reina milnovios;\\
terca ella sí, pero no fue capaz de guardarle al esposo\\
de ley el noble casal, hasta que se cumpliera el retorno».
\end{verse}

Quien no ha vivido en un pueblo pequeño no sabe lo malvada que llega a ser la gente, lo afiladas que pueden ser las lenguas. Penélope lleva veinte años viuda a todos los efectos y casi cuatro años resistiendo el asedio de un centenar de moscones, pero en cuanto se oye música, a las comadres les falta tiempo para condenarla. O puede que los más ávidos en la condena fueran los compadres, encantados de que por fin un pretendiente hubiera puesto a la indómita mujer en el lugar que le correspondía.

Odiseo también se pone guapo y, en lugar de colonia, Atenea le echa por encima la gracia de los inmortales. Hecho un brazo de mar, se sienta junto a su esposa y le dice:

\begin{verse}
«Bendita, a ti por encima de las femeniles mujeres\\
el corazón te secaron los que han en Olimpo su sede.\\
Ninguna otra mujer se estaría con ánimo inerte\\
sin acercarse al esposo que, luego de mil padeceres,\\
vuelve a su patria y su tierra tras veinte años ausente.\\
Ea ya, nana, prepárame un lecho para que me acueste:\\
de hierro, sí, el corazón que en el pecho guardado ella tiene».
\end{verse}

Nadie duda aquí de que picodeoro metisagaz se está haciendo el interesante y no lleva intención alguna de dormir solo. Penélope, sin embargo, no se achica:

\begin{verse}
«Hombre de dios, ni te estoy desdeñando ni soy altanera\\
ni estoy tan enmaravillada, que sé muy bien cómo eras\\
cuando de Ítaca echaste la nave larguirremera.\\
Así que, ea, aparéjale el lecho abrigado, Euriclea,\\
fuera del tálamo airoso que con sus manos hiciera;\\
sácale ahora aquí el sólido lecho y su cama adereza\\
con los tapices, la colcha y la lustrigalana estameña».
\end{verse}

A Penélope milpensamientos tampoco le faltan recursos ni astucia, pues la parrafada no es otra cosa que una treta. Odiseo responde ofendido (o quizás haciéndose el ofendido):

\begin{verse}
«Esta palabra que dices, mujer, en verdad me ha dolido;\\
¿quién a otra parte ha movido mi lecho? Aun a muy entendido\\
difícil, sí, le sería; a no ser que viniendo un dios mismo\\
con el su solo querer lo mudara sin pena a otro sitio.\\
Que de los hombres, ni aun mozo pujante, ningún mortal vivo\\
desencajarlo podría, pues guarda este lecho pulido\\
seña bien grande, que yo lo labré, y ninguno conmigo.\\
Dentro del patio creció piculanceolado un olivo,\\
alto en sazón y lozano; como una columna, fornido.\\
En torno a él levanté yo la alcoba, hasta dar cabo mismo,\\
con piedra bien apretada y la teché en buen abrigo,\\
y le hice sólidas puertas apareando sus quicios.

Luego corté la maraña de larga crin al olivo,\\
y, hecha la poda, desde la raíz dejé el tronco lampiño\\
a bronce en sabia labor y, adrizándolo a regla y a hilo,\\
en él labré el pie de cama calándolo con el arquillo.\\
Fui de ese pie cepillándome un lecho, hasta dar fin cumplido,\\
con oro obrando y con plata y con marfil sus aliños;\\
y una correa de toro tensé, de un granate encendido.\\
Esta es la seña que te manifiesto; y no sé si ahora mismo\\
sigue, mujer, allí firme este lecho, o ya en otro sitio\\
alguien lo puso, cortando por bajo ese tronco de olivo».
\end{verse}

Se trata del famoso olivo en el que está empotrada la cama de Odiseo y toda la escena es la clave del reconocimiento entre los esposos. Penélope, por fin, se deshiela y le suelta una parrafada memorable, en la que, por cierto, rompe una lanza a favor de Helena:

\begin{verse}
«No te amohínes conmigo, Odiseo, que siempre el más cuerdo\\
de entre los hombres tú has sido. Quebrantos nos ha dado el cielo,\\
harto celoso de que ambos gozáramos juntos el tiempo\\
de juventud y el umbral de la vida alcanzáramos viejos.\\
Así que no te me aíres ni lleves ahora en despecho\\
que al verte no te tuviera amor y reconocimiento.\\
Que siempre mi corazón recela y tiembla de miedo\\
por los que con sus palabras aquí me vengan mintiendo,\\
que en el urdir malas mañas se hicieron muchos maestros.\\
Helena Argiva tampoco ---y por padre es nacida de Zeus---\\
se habría unido en amor y delicias a un hombre extranjero\\
de haber sabido que luego los bravos hijos de aqueos\\
la llevarían de vuelta a su casa a la patria de nuevo.\\
Sin duda un dios la empujó a su malhacer deshonesto;\\
que ella antes nunca a ceguera había en su alma hecho hueco,\\
triste ceguera que para nosotros dolor fue primero.\\
Ahora, pues que de señas tan claras cuenta me has hecho\\
y es una cama la nuestra que nunca la vio moridero,\\
no, sino solo tú y yo y esa mi única ama que tengo,\\
Actóride, la que mi padre me dio cuando vine a este techo\\
y nos guardaba las puertas de nuestro cerrado aposento,\\
dejas rendido mi corazón, por más que era de hierro».
\end{verse}

Ulises no ha olvidado que le queda un viaje pendiente y así se lo dice a su esposa:

\begin{verse}
«Mujer, que no hemos llegado hasta el cabo final de las pruebas,\\
que aún habrá de venir por delante una inmensa faena,\\
larga y difícil, y a ella he de dar cumplimiento, que es fuerza.\\
Agüero fue que me hizo el ánima a mí de Tiresias\\
cuando a la casa de Hades en busca fui de respuesta\\
para el regreso de mis compañeros y el mío. Mas, ¡ea!,\\
ahora vayamos al lecho, mujer, y que juntos nos venga\\
del dulce sueño el regalo, uno del otro a la vera».
\end{verse}

Pero Penélope está desvelada y no le apetece meterse en el catre sin que Odiseo le revele el misterioso trabajo que todavía le queda pendiente. Al marido no le hace ni pizca de gracia seguir con la cháchara, pero se resigna y dice:

\begin{verse}
«Relato te haré, sin guardarme ni un verbo;\\
aunque alegría no vas a llevar, como yo no me alegro.\\
Mandó que plazas corriera de los mortales y pueblos\\
llevando un remo manobediente por mil derroteros\\
hasta que dé con un suelo de hombres que mar nunca vieron,\\
ni la comida mezclaron con sal, y que nunca supieron\\
nada de naves rostribermejas y nada de remos\\
bienllevaderos, que mismo son alas para los veleros.\\
De esto una seña muy clara me dio que ocultarte no quiero:\\
cuando me encuentre, rodando caminos, con otro andariego\\
que a mí me diga que al hombro gentil llevo palo de aviento,\\
que hunda me dijo él entonces el remo allí mismo en el suelo,\\
y a Poseidón soberano le inmole en aplacamiento\\
carnero, toro y verraco, y vuelva a esta casa de nuevo\\
para ofrecer una santa hecatombe a los no morideros,\\
los tantos dioses que pueblan la vasta campa del cielo,\\
uno tras otro hasta todos. Y a mí desde el mar salinero\\
muy fácil muerte me habrá de llegar, que el acabamiento\\
por apacible vejez abrumado vendrá; y con un pueblo\\
próspero siempre a mi lado. De todo auguró cumplimiento».
\end{verse}

\begin{verse}
Y dijo tras lo escuchado Penélope bienserenísima:\\
«Si a ti una buena vejez los dioses al fin te destinan\\
es que esperanza hay también de librarte de cuitas un día».
\end{verse}

Podemos dejar aquí, en un final feliz un poco falso, nuestra historia. No todos los cabos han quedado atados, sin embargo. Ulises le hará a su esposa el relato de sus aventuras y ella tendrá que fingir que no le importan el año con Circe ni los siete años con Calipso. O quizás más difícil, tendrá que aprender a convivir con un extraño.
